\subsection{\rev Pairwise comparisons of learning-signal betas per region}

\begin{table}[htbp]
    \centering
    \rev
    \begin{threeparttable}
    \caption{Paired $t$-tests between the three learning-signal parametric
    modulators of model7 (signed RPE, unsigned RPE (uRPE), Shannon surprise),
    computed on subject-level mean betas within each previously reported
    region ($n=58$). $p$-values are Holm-corrected across all
    108 pairwise tests.}
    \label{tab:roi_signal_ttests}
    \begin{tabularx}{\textwidth}{X rr rr rr}
        \toprule
        & \multicolumn{2}{c}{uRPE $-$ RPE} & \multicolumn{2}{c}{uRPE $-$ surprise} & \multicolumn{2}{c}{RPE $-$ surprise} \\
        \cmidrule(lr){2-3}\cmidrule(lr){4-5}\cmidrule(lr){6-7}
        \textbf{Region} & $t$ & $p$ & $t$ & $p$ & $t$ & $p$ \\
        \midrule
        \multicolumn{7}{l}{\textit{RPE regions}} \\
        Angular gyrus L & -8.56 & $<.001$ & -3.90 & $0.009$ & 7.59 & $<.001$ \\
        Caudate/putamen L & -15.04 & $<.001$ & -4.47 & $0.002$ & 14.91 & $<.001$ \\
        Caudate/putamen R & -15.18 & $<.001$ & -4.03 & $0.006$ & 15.62 & $<.001$ \\
        Cerebellum R & -8.30 & $<.001$ & -1.71 & $0.842$ & 9.57 & $<.001$ \\
        Visual cortex L (lingual) & -10.71 & $<.001$ & -3.43 & $0.032$ & 9.79 & $<.001$ \\
        Visual cortex L (mid. occ.) & -11.07 & $<.001$ & -5.44 & $<.001$ & 9.08 & $<.001$ \\
        vmPFC & -11.59 & $<.001$ & -6.62 & $<.001$ & 9.22 & $<.001$ \\
        \addlinespace
        \multicolumn{7}{l}{\textit{Surprise regions}} \\
        Angular gyrus L & -8.10 & $<.001$ & -6.33 & $<.001$ & 4.85 & $<.001$ \\
        Angular gyrus R & -6.67 & $<.001$ & -4.34 & $0.002$ & 4.35 & $0.002$ \\
        Precuneus L & -5.15 & $<.001$ & -4.46 & $0.002$ & 3.04 & $0.083$ \\
        dlPFC L & -7.70 & $<.001$ & -6.71 & $<.001$ & 4.76 & $<.001$ \\
        \addlinespace
        \multicolumn{7}{l}{\textit{uRPE regions}} \\
        Anterior insula R & 5.46 & $<.001$ & 6.08 & $<.001$ & -2.47 & $0.266$ \\
        Inferior parietal lobule R & -0.23 & $1.000$ & 4.48 & $0.002$ & 3.90 & $0.009$ \\
        Insula L & 2.94 & $0.098$ & 4.33 & $0.002$ & -0.20 & $1.000$ \\
        Middle frontal gyrus R & 2.44 & $0.266$ & 6.46 & $<.001$ & 1.62 & $0.888$ \\
        Orbital IFG R & 3.23 & $0.051$ & 4.43 & $0.002$ & -1.16 & $1.000$ \\
        Superior frontal sulcus R & 4.06 & $0.006$ & 4.45 & $0.002$ & -1.93 & $0.590$ \\
        dmPFC R & 4.35 & $0.002$ & 5.91 & $<.001$ & -0.96 & $1.000$ \\
        \addlinespace
        \bottomrule
    \end{tabularx}
    \begin{tablenotes}\footnotesize
    \item Betas are averaged over the six runs; all modulators were
    z-scored per run before entering the GLM, so betas share a common
    scale. Positive $t$: the first-named signal has the larger beta.
    \end{tablenotes}
    \end{threeparttable}
\end{table}%
